\section{Discussion}
\label{sec:discussion}
\subsection{Response fidelity is a separate model-selection criterion}
The central result is not that a particular loss or architecture wins every comparison. It is that Teacher-response fidelity, agreement with stated choices and execution consistency rank models differently. In the controlled benchmark, explicit response weighting improves finite-contrast error, but changing the choice specification produces a larger gain. The signed and direction+magnitude objectives are not clearly separated by the paired interval, and neither improves interaction error over soft KL. The evidence therefore supports evaluating and weighting responses, without establishing that the split sign--magnitude construction is essential.

This interpretation follows the information structure of the experiment. The same endpoint targets are available to every controlled model, and Eq.~\ref{eq:bound} connects their errors to response error. Paired supervision reallocates optimization effort across those errors rather than revealing additional Teacher behavior. Results across the six held-out personas and the available Teacher reaggregations support the within-benchmark comparison. Transfer to unseen intervention families, intensities or Teachers remains a separate empirical question.

MNL-S's advantage makes model specification an important control rather than a secondary baseline. Its linear utilities provide inspectable coefficients, and separate neural timing retains most of the timing accuracy of joint models. This combination supports modular estimation and later recalibration. It does not establish that travelers follow linear utilities or that every coefficient has an economically valid interpretation: the choice target is an LLM judgment, and smoothness or regularization could also explain the gain. Behavioral interpretation requires checking identified utility differences, relevant signs and implied sensitivities, not merely naming the model logit.

\subsection{Human agreement requires an independent empirical test}
Teacher fidelity is not a proxy for agreement with people. MNL-S reverses the stated Singapore accessibility response, whereas CE+KL, the least faithful controlled objective, comes closest to that contrast. Baseline levels and responses also separate: SA-Student understates Singapore PT use by about 31 percentage points yet has the smallest mean absolute response bias there; MNL-S comes within four points of the PT share but misses important changes. These observations extend the distinction between predictive and behavioral criteria \citep{wang2020deep,martinbaos2023prediction} to models compressed from LLM judgments.

The contemporary API reference helps identify where follow-up tests are needed, without assigning historical causation. For Singapore accessibility, it also attenuates the stated response, so bringing SA-Student closer to that reference would leave substantial human disagreement. For Shanghai delay under the reference encoding, the current model is close to the human point estimate while SA-Student responds much more strongly. The second pattern motivates a controlled check of Student specification and input encoding; it does not establish that compression caused the discrepancy. The Shanghai walking-versus-waiting result additionally illustrates cancellation: a small total bias can combine larger opposing Student--reference and reference--human terms.

This distinction suggests two different development tasks. Compression should be evaluated against frozen Teacher targets. Empirical alignment should be fitted and selected using separate human data, then tested on people and tasks not used for that alignment, as motivated by recent travel-choice and travel-satisfaction research \citep{liu2026persona,xu2026satisfaction}. The surveys in the present study were reserved for evaluation, so the observed differences diagnose calibration needs rather than demonstrate a calibrated traveler model.

\subsection{Routability, response transmission and practical use}
The Helsinki experiment isolates what happens after behavioral prediction. For SA-Student, baseline PT use falls by more than eleven points between prediction and boarding while the deterministic delay response changes by less than one point. A small response gap can thus coexist with a substantial level gap. Feasibility-constrained assignment removes PT choices for which the specified search finds no connection, yet does not reduce the response gap for this population. Routability and preservation of a prescribed behavioral distribution are different objectives.

The stagewise accounting makes these differences auditable. All shares use the same initial population; predicted and adjusted probabilities remain distinct from assigned modes, routed modes and boarding events. Signed and absolute gaps are both needed because cancellation across fitted models can conceal errors. Timing outputs also warrant separate evaluation: they change itineraries even when binary feasibility remains unchanged. The additional speed experiment shows why a stable mode count alone cannot validate a journey-time outcome.

For implementation, the relevant efficiency result concerns repeated use after training. Behavioral inference consumes approximately 0.35\% of the measured larger scenario builds, shifting attention to supply calculation and reuse rather than marginal differences between compact Students. The retained token ledger quantifies an additional acquisition-cost component, but it does not yet determine the lifecycle break-even workload: unmetered attempts and local-compute costs remain outside the estimate in Table~\ref{tab:api_cost}.

\subsection{Scope of inference and next tests}
Three boundaries define the conclusions. First, the compression comparison uses one archived Teacher and six held-out synthetic personas. Teacher reaggregation and seed variation test sensitivity within that design; they do not substitute for new personas, unseen interventions or another frozen Teacher. Second, the human evidence concerns stated responses in two convenience samples under documented encodings. The cities used different task formats, Shanghai's qualitative tasks require numerical reference values, and the existing sensitivity analysis holds its delay at 15 minutes. Explicit currency, alternative-set interpretation and randomized, numerically matched tasks are priorities for further validation. Third, Helsinki holds physical supply and behavioral predictions fixed to isolate execution effects. Physical service changes and feedback from experienced conditions require separate experiments before extending the conclusions to equilibrium or operational disruption forecasting.

These boundaries define a usable validation sequence rather than a single pass/fail imitation score: specify the intervention and output; test compression against archived targets; compare the corresponding human response; and audit the outcome after assignment and execution. A failure at one level identifies the next model component or evidence source to examine. Supplementary Section~\ref{supp-sec:materials_scope} provides the versioned material inventory and distinguishes reproducible target-based evaluation from live-API and restricted-data dependencies.
